\documentclass[11pt,letterpaper]{article}

% --- Packages ---
\usepackage[utf8]{inputenc}
\usepackage[T1]{fontenc}
\usepackage{lmodern}
\usepackage[margin=1in]{geometry}
\usepackage{amsmath,amssymb}
\usepackage{booktabs}
\usepackage{graphicx}
\usepackage{hyperref}
\usepackage{xcolor}
\usepackage{enumitem}
\usepackage{tabularx}
\usepackage{caption}
\usepackage{float}
\usepackage{microtype}
\usepackage{natbib}
\usepackage{authblk}

\hypersetup{
  colorlinks=true,
  linkcolor=blue!70!black,
  citecolor=blue!70!black,
  urlcolor=blue!70!black
}

% --- Title ---
\title{\textbf{MergeConflictBench: Evaluating LLM Agents on\\Semantically Correct Merge Conflict Resolution}}

\author[1]{Daniel Chen}
\affil[1]{Anything, Inc.}
\date{}

\begin{document}
\maketitle

% --- Abstract ---
\begin{abstract}
We introduce \textsc{MergeConflictBench}, an evaluation suite for measuring the ability of LLM agents to resolve real merge conflicts while preserving behavior from both branches. The benchmark consists of 86 cases extracted from production merge conflict tracking in a commercial React-family codebase---React.js and Next.js web surfaces, React Native mobile surfaces, and adjacent JavaScript/JSX API and helper modules. The corpus spans 288 conflicted files with 1{,}063 conflict blocks across domains including authentication, mobile UI, regulatory data processing, animation systems, fleet management, e-commerce, and CRM tooling. Each case is paired with a hidden holdout test suite---2{,}540 executable tests in total across all 86 cases---that verifies behaviors originating from all three code lineages: the common base, the ``ours'' branch, and the ``theirs'' branch. An agent's task is to produce resolved files free of conflict markers; its output is scored by executing the hidden tests, directly measuring whether the resolution correctly integrates both branches' contributions. This design extends the principle established by \textsc{RefactorBench}---\emph{behavioral preservation is a functional property best verified by functional tests}---from single-branch refactoring to multi-branch integration, where the agent must synthesize semantically compatible code from two divergent change sets. Release-time validation executes all 2{,}540 hidden tests against the checked-in reference resolutions, with all 86 cases passing; deterministic accept-ours, accept-theirs, and accept-both baselines pass 17, 10, and 48 cases respectively. A one-shot LLM block-replacement harness across five contemporary models passes between 25 and 36 cases, showing that semantic merge resolution remains challenging even when models can inspect the full conflicted file contents.
\end{abstract}

\medskip
\noindent\textbf{Keywords:} benchmark, React.js, merge conflict resolution, large language models, agentic evaluation, behavioral preservation

% ============================================================
\section{Introduction}
\label{sec:introduction}
% ============================================================

Merge conflicts are an unavoidable consequence of collaborative software development. When two branches modify overlapping regions of the same file, version control systems cannot automatically reconcile the changes and instead produce \emph{conflict markers}---the familiar \texttt{<{}<{}<{}<{}<{}<{}<}, \texttt{=======}, and \texttt{>{}>{}>{}>{}>{}>{}>{}} delimiters that bracket the competing versions. A developer must then manually inspect each conflict block and produce a resolution that correctly integrates the intent of both branches.

For trivial conflicts---a renamed variable on one side, an added import on the other---resolution is mechanical. But many real-world conflicts require \emph{semantic understanding}: when both branches restructure the same function, add interacting features, or change overlapping state management logic, the correct resolution is not a simple union of the two sides but a careful synthesis that preserves the behavior introduced by each branch while maintaining overall program coherence.

Three common resolution strategies illustrate the difficulty spectrum:

\begin{itemize}[leftmargin=2em]
\item \textbf{Accept ours / accept theirs.} Discard one branch's changes entirely. Correct only when one side's changes are strictly superseded.
\item \textbf{Accept both.} Concatenate both sides. Correct only when the changes are independent and order-insensitive.
\item \textbf{Custom resolution.} The developer writes new code that integrates both branches' intent. This is the non-trivial case.
\end{itemize}

LLM-based agents can potentially resolve conflicts that require semantic understanding, but evaluating whether an agent's resolution is correct poses the same challenge that motivated \textsc{RefactorBench}~\citep{chen2025refactorbench}: structural comparison (exact match, diff size) is unreliable for code transformations where multiple textually distinct outputs can be functionally equivalent. A resolution that reorders function parameters, uses a different variable name, or restructures a conditional may be perfectly correct despite differing from the reference.

We address this with hidden functional tests. Each benchmark case includes a test suite---invisible to the resolving agent---that exercises behaviors from all three code lineages: the common ancestor (base), the current branch (ours), and the incoming branch (theirs). A resolution passes if and only if it preserves all three categories of behavior.

This paper introduces \textsc{MergeConflictBench}, a benchmark of 86 React-family merge conflict cases from production, and makes the following contributions:

\begin{enumerate}[leftmargin=2em]
\item \textbf{\textsc{MergeConflictBench} itself}: a benchmark of 86 real React.js/Next.js/React Native merge conflict cases spanning 288 files and 1{,}063 conflict blocks, with 2{,}540 executable hidden tests covering behaviors from base, ours, and theirs branches across all cases.

\item \textbf{Extension of behavioral evaluation to merge resolution}: we demonstrate that the hidden-test methodology from \textsc{RefactorBench} transfers naturally to the merge conflict domain, where it must verify the \emph{synthesis} of two change sets rather than the \emph{preservation} of one.

\item \textbf{Computed release-time and model baselines}: we report a full reference-resolution validation, three deterministic merge-strategy baselines, and a five-model one-shot LLM matrix, with all counts generated from checked-in benchmark artifacts or the released Escher eval harness.
\end{enumerate}

% ============================================================
\section{Related Work}
\label{sec:related}
% ============================================================

\paragraph{Behavioral evaluation of code transformations.} \textsc{RefactorBench}~\citep{chen2025refactorbench} introduced hidden holdout test suites for evaluating behavior-preserving code decomposition, demonstrating that functional tests detect regressions invisible to compilation, code similarity, and LLM-as-judge metrics. \textsc{MergeConflictBench} extends this methodology from single-branch transformations (refactoring) to multi-branch integration (merge conflict resolution), where the correctness criterion is not ``preserve existing behavior'' but ``synthesize the behaviors of two divergent branches.''

\paragraph{Agentic coding benchmarks.} SWE-bench~\citep{jimenez2024swebench} evaluates agents on resolving real GitHub issues, using repository test suites as ground truth. SWE-agent~\citep{yang2024sweagent} demonstrated that agent tool design significantly affects performance on SWE-bench. OpenHands~\citep{wang2024openhands} provides an open platform for coding agents. Aider's code editing benchmarks~\citep{gauthier2024aider,gauthier2024polyglot} use test-based evaluation for refactoring and editing tasks across languages. These benchmarks evaluate \emph{generative} coding tasks (write new code to fix an issue); \textsc{MergeConflictBench} evaluates an \emph{integrative} task (combine two existing change sets correctly).

\paragraph{Merge conflict resolution.} Traditional merge tools (kdiff3, meld, IntelliJ's merge resolver) provide three-way diff views but leave resolution to the developer. DeepMerge~\citep{dinella2022deepmerge} trains sequence-to-sequence models on merge resolution data, evaluating by exact match against reference resolutions. MergeBERT~\citep{svyatkovskiy2022mergebert} uses a token-level classification approach. MergeGen~\citep{zhang2024mergegen} applies retrieval-augmented generation. These approaches evaluate by \emph{textual similarity} to reference resolutions---a metric that can over-penalize correct resolutions that differ from the reference and cannot verify semantic correctness. \textsc{MergeConflictBench} is, to our knowledge, the first merge conflict benchmark that evaluates by \emph{executing behavioral tests} against the resolution.

\paragraph{Semantic merge tools.} Semantic merge approaches~\citep{apel2011semistructured} use language-aware AST merging to reduce spurious conflicts but do not handle semantic conflicts where both branches modify the same logic. Program synthesis approaches to merge~\citep{sousa2018verified} can verify correctness via formal specifications but are limited to restricted program fragments. \textsc{MergeConflictBench} targets the general case where arbitrary code changes must be integrated.

% ============================================================
\section{Benchmark Design}
\label{sec:design}
% ============================================================

\subsection{Design Principles}

\textsc{MergeConflictBench} is built on four design principles:

\begin{enumerate}[leftmargin=2em]
\item \textbf{Functional evaluation.} The primary metric is whether hidden holdout tests pass after resolution, directly measuring whether the agent correctly integrated both branches' contributions. Structural metrics (no remaining conflict markers, compilation success) are necessary conditions recorded as secondary scores.

\item \textbf{Three-lineage behavioral coverage.} The hidden test suite for each case covers behaviors from all three code lineages---base, ours, and theirs. A resolution that simply accepts one side will fail the tests for the other side's behaviors. A resolution that accepts both without proper integration may fail tests that exercise interactions between the two branches' changes.

\item \textbf{Real production conflicts.} All cases are extracted from production merge conflict tracking, not synthesized or curated from open-source repositories. This ensures the conflicts represent the kinds of semantic entanglement that occur in real collaborative development.

\item \textbf{Non-trivial resolution required.} Cases are selected for conflicts that were resolved via \texttt{custom\_resolution} or \texttt{accepted\_both} strategies in production---not \texttt{accepted\_ours} or \texttt{accepted\_theirs}, which are trivial.
\end{enumerate}

\subsection{Data Source}

Cases are extracted from production merge conflict tracking stored in BigQuery. The source tables contain 3.3 million conflict tracking rows (\texttt{generation.fs\_conflict\_results}) and 109{,}000 resolution tracking rows (\texttt{generation.merge\_fs\_conflict\_resolution\_results}). The extraction pipeline selects resolved conflicts where the resolution strategy was \texttt{custom\_resolution} or \texttt{accepted\_both}, filters for cases with at least two conflict blocks (to exclude trivially small conflicts), and retrieves the conflicted files (with markers) and the reference resolution.

\subsection{Three-Branch Structure}

Each benchmark case captures the three-way merge structure:

\begin{itemize}[leftmargin=2em]
\item \textbf{Base}: the common ancestor before either branch diverged.
\item \textbf{Ours}: the current branch's version (changes since base).
\item \textbf{Theirs}: the incoming branch's version (changes since base).
\item \textbf{Conflicted}: the merge result with conflict markers, containing both \texttt{ours} and \texttt{theirs} content delimited by \texttt{<{}<{}<{}<{}<{}<{}<}, \texttt{=======}, and \texttt{>{}>{}>{}>{}>{}>{}>{}} markers.
\end{itemize}

The agent receives only the conflicted files. The base, ours, and theirs versions are used during test authoring but are not provided to the agent at evaluation time. The reference resolution is used as a correctness guide during test authoring but is not part of the scoring---only the hidden tests determine correctness.

\subsection{Case Layout}

Each case directory follows a consistent structure:

\begin{verbatim}
  data/<case_name>/
    eval.config.json              # case metadata
    conflict_eval.config.json     # conflict blocks + test config
    conflicted/                   # files with markers (agent input)
    resolved/                     # reference resolution (test authoring)
    merge.test.js                 # hidden test suite (holdout)
\end{verbatim}

The \texttt{eval.config.json} records the case name, description, source conflict ID, timestamp, language, number of conflicted files, and total conflict block count. The \texttt{conflict\_eval.config.json} enumerates every conflict block (with ours/theirs content), lists the conflicted file paths, specifies the test file, and documents the preserved behaviors as a structured contract.

\subsection{Hidden Test Suite Design}

Each hidden test suite is authored by a senior engineer following these principles:

\begin{itemize}[leftmargin=2em]
\item \textbf{Test observable behavior from all three lineages.} Test suites are authored against behaviors identified from the base, ours, and theirs lineages, and some behavior-contract entries include explicit origin tags where that attribution was recorded.

\item \textbf{Cover the preserved-behavior contract.} The \texttt{preservedBehaviors} field in each case's config serves as the contract: a list of human-readable behavior descriptions. The test suite is the executable form of this contract.

\item \textbf{Use stable assertions.} Test by return values, API responses, rendered output, and state transitions---not by internal variable names or code structure, which may differ in a correct alternative resolution.

\item \textbf{Exercise interactions.} Where the two branches' changes interact (e.g., both modify the same data structure, both add validation rules to the same endpoint), include tests that verify the combined behavior, not just each branch's changes in isolation.
\end{itemize}

\subsection{Scoring}

\textsc{MergeConflictBench} produces five scores per resolution attempt:

\begin{table}[ht]
\centering
\caption{Scoring metrics for \textsc{MergeConflictBench}.}
\label{tab:scores}
\small
\begin{tabularx}{\textwidth}{llX}
\toprule
\textbf{Score} & \textbf{Type} & \textbf{Description} \\
\midrule
Passes Tests & Binary & Do the hidden holdout tests pass against the resolved files? \emph{Primary metric.} \\
Resolution Reported Success & Binary & Did the resolver produce a complete marker-free candidate? \\
Block Resolution Coverage & Ratio & Fraction of configured conflict blocks for which the resolver produced a parseable replacement. \\
No Conflict Markers & Binary & Output contains no remaining \texttt{<{}<{}<{}<{}<{}<{}<} / \texttt{=======} / \texttt{>{}>{}>{}>{}>{}>{}>{}} markers. \\
Usage Available & Binary & Did the provider route report usage metadata for the model call? \\
\bottomrule
\end{tabularx}
\end{table}

The \emph{Passes Tests} score is the primary metric. The other scores provide diagnostic information: a resolution that passes tests but has remaining conflict markers would indicate a test suite gap; partial block coverage distinguishes malformed outputs from behavior-preservation failures; a resolution that the resolver reports as successful but that fails tests reveals miscalibration.

% ============================================================
\section{Corpus}
\label{sec:corpus}
% ============================================================

The corpus consists of 86 cases totaling 288 conflicted files, 1{,}063 conflict blocks, and 2{,}540 executable hidden tests. All cases are JavaScript/JSX drawn from React-family application code: React.js and Next.js web surfaces, React Native mobile surfaces, and adjacent API/helper modules that support those applications. Block counts range from 2 to 46, file counts from 1 to 8, and per-case test counts from 9 to 94. Counts are generated from the checked-in configs and reference-validation report using \texttt{scripts/summarize-corpus.mjs}. Table~\ref{tab:corpus} presents a representative sample of 40 cases spanning all complexity tiers, followed by aggregate statistics for the full corpus.

\begin{table}[H]
\centering
\caption{\textsc{MergeConflictBench} corpus (representative sample). Files = number of conflicted files. Blocks = total conflict blocks. Tests = executable hidden holdout tests per case. Domain describes the functional area of the code.}
\label{tab:corpus}
\small
\begin{tabular}{llrrrl}
\toprule
\textbf{Case} & \textbf{Tier} & \textbf{Files} & \textbf{Blocks} & \textbf{Tests} & \textbf{Domain} \\
\midrule
\texttt{fleet\_parts\_transfers}   & Complex  & 4 & 46 & 44  & Fleet parts transfer workflow \\
\texttt{state\_parser\_apply}      & Complex  & 1 & 45 & 40  & Regulatory data import API \\
\texttt{account\_statements}       & Complex  & 2 & 36 & 35  & Account statement generation \\
\texttt{fleet\_parts\_receipts}    & Complex  & 3 & 34 & 34  & Fleet parts receipt processing \\
\texttt{stats\_home\_page}         & Complex  & 2 & 33 & 29  & Statistics dashboard home \\
\texttt{mobile\_auth\_flow}        & Complex  & 4 & 32 & 55  & Mobile auth via WebView \\
\texttt{dispatch\_sidebar\_nav}    & Complex  & 2 & 31 & 28  & Dispatch sidebar navigation \\
\texttt{account\_auth\_pages}      & Complex  & 4 & 30 & 63  & Signup/signin/reset pages \\
\texttt{state\_parser\_preview}    & Complex  & 8 & 25 & 94  & Admin preview with registry \\
\texttt{campaign\_mapping\_tool}   & Complex  & 7 & 25 & 63  & Campaign canvas \& toolbar \\
\texttt{listing\_watchlist}        & Complex  & 3 & 25 & 22  & Listing watchlist feature \\
\texttt{celebration\_overlay}      & Complex  & 1 & 24 & 72  & Confetti animation \& audio \\
\texttt{mobile\_chat\_profile}     & Complex  & 5 & 22 & 20  & Mobile chat profile views \\
\texttt{bulk\_product\_status}     & Complex  & 3 & 22 & 32  & Bulk status change \& preview \\
\texttt{legal\_compliance\_settings}& Complex & 2 & 22 & 32  & EULA/SMS document URLs \\
\texttt{mobile\_task\_detail}      & Complex  & 5 & 21 & 62  & Task modal with comments \\
\texttt{route\_mileage\_planning}  & Complex  & 3 & 21 & 34  & Mileage estimation hooks \\
\texttt{navigation\_layout\_tailwind}& Complex& 3 & 20 & 31  & Navigation layout styling \\
\midrule
\texttt{resource\_crud\_modals}    & Moderate & 7 & 19 & 29  & Resource CRUD modal forms \\
\texttt{checkout\_payment\_verify} & Moderate & 3 & 19 & 28  & Checkout payment verification \\
\texttt{trade\_show\_form}         & Moderate & 1 & 19 & 83  & Mobile trade show creation \\
\texttt{dealer\_integration\_csv}  & Moderate & 3 & 17 & 27  & Dealer CSV integration \\
\texttt{onboarding\_agreement}     & Moderate & 4 & 16 & 60  & PDF/HTML agreement flow \\
\texttt{invoice\_create\_modal}    & Moderate & 4 & 15 & 46  & Invoice creation modal \\
\texttt{mobile\_tab\_navigation}   & Moderate & 4 & 14 & 45  & Bottom tab bar \& badges \\
\texttt{ai\_studio\_animation}     & Moderate & 2 & 13 & 40  & AI studio animations \\
\texttt{org\_switcher}             & Moderate & 3 & 12 & 60  & Org switcher with RBAC \\
\texttt{legal\_compliance\_v2}     & Moderate & 2 & 12 & 25  & Document link storage \\
\texttt{globe\_renderer\_animation}& Moderate & 4 & 10 & 33  & Globe renderer animation \\
\texttt{auth\_api\_routes}         & Moderate & 4 &  8 & 52  & Login/signup/OTP routes \\
\midrule
\texttt{billing\_pricing\_pages}   & Simple   & 3 &  7 & 18  & Billing \& pricing pages \\
\texttt{athlete\_contact\_api}     & Simple   & 1 &  6 & 30  & Contact endpoint \& email rules \\
\texttt{signin\_academy\_service}  & Simple   & 2 &  6 & 22  & Signin \& academy service \\
\texttt{password\_reset\_full}     & Simple   & 5 &  5 & 22  & Password reset flow (full) \\
\texttt{mobile\_customer\_filters} & Simple   & 4 &  4 & 47  & Customer search \& team queries \\
\texttt{presales\_copilot\_products}& Simple  & 4 &  4 & 39  & Presales copilot products \\
\texttt{campaign\_call\_panel}     & Simple   & 1 &  4 & 31  & Call panel \& speaker tracking \\
\texttt{sme\_leader\_domains}      & Simple   & 3 &  3 & 18  & SME leader domains \\
\texttt{product\_import\_modal}    & Simple   & 1 &  2 & 50  & Field mapping for imports \\
\texttt{matchmaking\_user\_helpers}& Simple   & 2 &  2 & 12  & Matchmaking user helpers \\
\midrule
\multicolumn{2}{l}{\emph{Remaining 46 cases}} & \multicolumn{4}{l}{\emph{(see supplementary data for full inventory)}} \\
\midrule
\textbf{Full corpus (86 cases)} &          & \textbf{288} & \textbf{1{,}063} & \textbf{2{,}540} & \\
\bottomrule
\end{tabular}
\end{table}

\subsection{Complexity Distribution}

Cases are stratified into three complexity tiers based on the number of conflict blocks:

\begin{itemize}[leftmargin=2em]
\item \textbf{Complex (20+ blocks):} 18 cases (21\%), containing 514 conflict blocks across 62 files with 790 hidden tests. These involve extensive interleaving of changes, often requiring the agent to understand how both branches restructured the same functions, state management, or API surfaces. The most complex case (\texttt{fleet\_parts\_transfers}) has 46 conflict blocks across 4 files---a fleet parts transfer workflow where both branches restructured transfer logic, receipt handling, and validation. The single-file maximum is \texttt{state\_parser\_apply} with 45 blocks in one file---a regulatory data import API where both branches restructured function signatures, parameter handling, and handler logic.

\item \textbf{Moderate (8--19 blocks):} 29 cases (34\%), containing 370 conflict blocks across 114 files with 956 hidden tests. These involve meaningful conflicts across multiple files but with more localized changes. For example, \texttt{onboarding\_agreement} (16 blocks across 4 files) involves conflicts in invitation flow state building, parameter passing, and rendering mode (iframe vs.\ direct). The median case in this tier has 13 blocks across 3 files.

\item \textbf{Simple (2--7 blocks):} 39 cases (45\%), containing 179 conflict blocks across 112 files with 794 hidden tests. Despite having fewer conflict blocks, these cases still require semantic understanding beyond accept-ours/theirs. For example, \texttt{athlete\_contact\_api} (6 blocks in 1 file) involves conflicting changes to email uniqueness rules, contact type checking, and error response formatting. Several simple cases have high test counts (e.g., \texttt{mobile\_customer\_filters} with 47 tests for 4 blocks), reflecting dense behavioral coverage even in structurally small conflicts.
\end{itemize}

The overall distribution has a median of 8.5 blocks per case (mean 12.4), with the bulk of cases in the simple-to-moderate range. This reflects the natural distribution of merge conflicts in production: most conflicts are localized, but a meaningful tail of highly complex multi-block conflicts exists and poses the greatest challenge for automated resolution.

\subsection{Conflict Type Taxonomy}

Across the corpus, we observe recurring categories of semantic conflict:

\begin{itemize}[leftmargin=2em]
\item \textbf{Function refactoring conflicts.} Both branches restructure the same function: one extracts parameters into an options object, the other adds new parameters. The resolution must combine the structural change with the new functionality.

\item \textbf{State management conflicts.} Both branches modify how component state is managed: one renames state variables, the other adds new state fields. The resolution must apply the renaming to both old and new fields.

\item \textbf{Import/dependency conflicts.} Both branches add imports from different modules or restructure existing imports. Typically the simplest conflict type, but can interact with other conflicts if the imported symbols are used in conflicting code.

\item \textbf{UI layout conflicts.} Both branches modify the same component's JSX structure: one changes the layout hierarchy, the other adds new UI elements. The resolution must place the new elements within the revised hierarchy.

\item \textbf{API contract conflicts.} Both branches change the same API endpoint: one modifies the request format, the other changes the response format. The resolution must be internally consistent across request and response handling.
\end{itemize}

\subsection{Preserved Behavior Contract}

Each case includes a \texttt{preservedBehaviors} list documenting what the hidden tests verify. Across the released corpus, these configs contain 1{,}383 human-readable behavior descriptions. Most entries are description-only; entries may also include an optional \texttt{origin} tag (base, ours, or theirs) when attribution has been explicitly recorded. For example, the \texttt{athlete\_contact\_api} case has 19 origin-tagged preserved behaviors: 10 from base (existing validation, auth checks, constraint error handling), 4 from ours (new email conflict detection logic), and 5 from theirs (contact type checking, updated error responses). This structure makes coverage gaps visible during test authoring while the executable test suite remains the scoring authority.

% ============================================================
\section{Evaluation Methodology}
\label{sec:methodology}
% ============================================================

\subsection{Agent Task}

The agent is given the conflicted files---source files containing \texttt{<{}<{}<{}<{}<{}<{}<} / \texttt{=======} / \texttt{>{}>{}>{}>{}>{}>{}>{}} conflict markers---and asked to produce resolved files. The agent may use any strategy: reading the conflict markers to understand both sides, reasoning about the semantic intent of each branch, and writing resolved code that integrates both contributions. The agent does \emph{not} have access to the hidden test suite, the reference resolution, or the base/ours/theirs versions as separate files.

\subsection{Evaluation Pipeline}

After the agent produces its resolution, the evaluation pipeline (\texttt{scripts/evaluate.ts}) performs the following steps:

\begin{enumerate}[leftmargin=2em]
\item Copy the agent's resolved files into a temporary workspace.
\item Copy the hidden test file (\texttt{merge.test.js}) into the workspace.
\item Execute the test suite via Vitest.
\item Record which tests pass and which fail.
\item Check for remaining conflict markers in the resolved files.
\item Report the five scoring dimensions (Table~\ref{tab:scores}).
\end{enumerate}

\subsection{Scoring Dimensions}

The primary metric is \emph{Passes Tests}: a binary indicator of whether all hidden holdout tests pass. This directly measures functional correctness---whether the resolution preserves behaviors from all three lineages.

Secondary metrics provide diagnostic context:

\begin{itemize}[leftmargin=2em]
\item \emph{No Conflict Markers} detects resolutions where the agent failed to address some conflict blocks, leaving raw markers in the output.
\item \emph{Block Resolution Coverage} measures whether the resolver supplied parseable replacements for the configured conflict blocks.
\item \emph{Resolution Reported Success} measures resolver calibration---the divergence between a complete marker-free candidate and actual functional correctness.
\item \emph{Usage Available} records whether the provider route returned usage metadata, which helps diagnose infrastructure failures separately from semantic failures.
\end{itemize}

% ============================================================
\section{Release-Time Validation, Baselines, and Model Matrix}
\label{sec:results}
% ============================================================

We compute two release-time checks from the public repository. First, we validate the checked-in reference resolutions by running every hidden test suite against each case's \texttt{resolved/} directory. Second, we evaluate three deterministic merge-editor baselines: accept ours, accept theirs, and accept both. These baselines replace every conflict block mechanically and do not perform semantic integration. We additionally report a one-shot LLM matrix using a tiny block-replacement harness in Escher: each model receives the conflicted files only, emits replacement text for each conflict block, and is scored by the same hidden tests.

All results in this section are generated by the commands documented in \path{data/eval-results/README.md}. The case and tier summaries are generated by \texttt{scripts/summarize-corpus.mjs}; the deterministic baselines are generated by \texttt{scripts/evaluate-naive-baselines.mjs}; and the model matrix is generated by the Escher merge-conflict eval harness. The raw Vitest JSON reports and CSV summaries are released under \path{data/eval-results/}.

\begin{table}[ht]
\centering
\caption{Release-time validation and deterministic baselines on all 86 cases. ``Cases Passed'' is the primary metric: a case passes only when its full hidden test suite passes. Assertion counts for failed deterministic baselines are lower bounds because syntax or import failures can abort a suite before every assertion runs.}
\label{tab:release-baselines}
\small
\begin{tabularx}{\textwidth}{lrrrX}
\toprule
\textbf{Resolution} & \textbf{Cases Passed} & \textbf{Case Rate} & \textbf{Assertions Passed} & \textbf{Description} \\
\midrule
Reference \texttt{resolved/} & 86 / 86 & 100.0\% & 2{,}540 / 2{,}540 & Checked-in production resolution used for test authoring \\
Accept ours & 17 / 86 & 19.8\% & 2{,}121 before failures & Use only the current-branch side of every conflict block \\
Accept theirs & 10 / 86 & 11.6\% & 2{,}025 before failures & Use only the incoming-branch side of every conflict block \\
Accept both & 48 / 86 & 55.8\% & 2{,}032 before failures & Concatenate ours followed by theirs for every conflict block \\
\bottomrule
\end{tabularx}
\end{table}

\begin{table}[ht]
\centering
\caption{One-shot LLM block-replacement results on all 86 cases. ``Cases Passed'' counts missing or unscored rows as failures, because the benchmark metric is performance across the full fixture set. ``Blocks Covered'' is mean parseable block-replacement coverage; ``Marker-free'' is the fraction of rows with no remaining conflict markers.}
\label{tab:model-matrix}
\small
\begin{tabular}{lrrrrr}
\toprule
\textbf{Model} & \textbf{Cases Passed} & \textbf{Case Rate} & \textbf{Blocks Covered} & \textbf{Marker-free} & \textbf{Usage} \\
\midrule
Claude Opus 4.6 & 36 / 86 & 41.9\% & 87.0\% & 75.6\% & 76 / 86 \\
Gemini 2.5 Pro & 35 / 86 & 40.7\% & 35.5\% & 27.9\% & 86 / 86 \\
Claude Sonnet 4.6 & 30 / 86 & 34.9\% & 84.5\% & 61.6\% & 76 / 86 \\
Gemini 2.5 Flash & 29 / 86 & 33.7\% & 65.3\% & 38.4\% & 86 / 86 \\
GPT-4.1 & 25 / 86 & 29.1\% & 84.0\% & 41.9\% & 86 / 86 \\
\bottomrule
\end{tabular}
\end{table}

\subsection{Key Observations}

\paragraph{The released reference resolutions are internally validated.} All 86 reference resolutions pass their hidden suites, for 2{,}540 passing executable tests and zero failures. This does not measure agent performance; it validates that the committed reference resolutions and hidden tests are mutually consistent at release time.

\paragraph{Naive merge-editor strategies are insufficient.} Accepting one side performs poorly: accept-ours passes 17 cases (19.8\%), while accept-theirs passes 10 cases (11.6\%). Accept-both is stronger, passing 48 cases (55.8\%), but still fails 38 cases. This supports the benchmark's selection criterion: many cases require semantic synthesis rather than a blanket policy.

\paragraph{Case-level pass rate is the appropriate baseline metric.} Failed baseline runs may abort early because the mechanically resolved file has syntax, import, or runtime errors. For this reason, assertion counts are diagnostic only; a case is counted as correct only when the full hidden suite passes.

\paragraph{One-shot LLM resolution is not yet enough.} The strongest model in the one-shot block-replacement harness, Claude Opus 4.6, passes 36 of 86 cases (41.9\%), below the deterministic accept-both baseline's 48 cases. This does not imply that LLM agents are worse than naive merging in general: the harness intentionally provides no tools, no test feedback, and no iterative repair loop. It does show that semantic merge resolution is not solved by a single model call that merely replaces conflict blocks.

\paragraph{Diagnostic scores reveal different failure modes.} Claude Opus and Sonnet produce high block-replacement coverage but lose cases to provider-route failures and hidden-test mismatches. Gemini 2.5 Pro passes 35 cases despite low marker-free and block-coverage diagnostics, indicating that some test suites do not exercise every unresolved region and that diagnostic scores should be reported alongside the primary hidden-test metric. The direct GPT-4.1 run has complete usage coverage and no provider-route failures, but lower semantic pass rate.

% ============================================================
\section{Discussion}
\label{sec:discussion}
% ============================================================

\subsection{Challenges in Test Authoring}

Writing hidden test suites for merge conflict cases poses challenges distinct from refactoring test authoring:

\begin{itemize}[leftmargin=2em]
\item \textbf{Three-lineage coverage.} The test author must understand the behavior of three code versions (base, ours, theirs) and verify that all three lineages' behaviors are preserved in the resolution. This requires careful reading of diffs to identify what each branch introduced.

\item \textbf{Interaction testing.} When both branches modify the same data structure or API, the test must verify the \emph{combined} behavior, not just each branch's changes independently. For example, if ours adds email validation and theirs adds phone validation to the same contact endpoint, a test must verify that both validations are active simultaneously.

\item \textbf{Stable assertions against divergent implementations.} Because multiple textually distinct resolutions can be correct, tests must assert on externally observable behavior (API responses, rendered output, return values) rather than internal structure. This is particularly challenging for UI component conflicts where the correct resolution might have different component hierarchies.
\end{itemize}

\subsection{Interpreting the Baselines}

The deterministic baselines are intentionally weak but informative. Accept-ours and accept-theirs measure how often one side subsumes the other; both rates are low. Accept-both measures how often textual concatenation is sufficient; it passes a slim majority of cases but still fails 44.2\% of the corpus. These failures include cases where concatenation creates syntax errors, duplicate definitions, inconsistent API contracts, or behavior that compiles but drops an interaction between the two branches.

The deterministic baseline results establish a floor for non-semantic merge policies and validate that many cases require integration decisions that are sensitive to behavior, not just conflict-marker removal. The one-shot LLM results establish a complementary baseline for tool-free model reasoning: models can often produce parseable replacements, but without iterative execution feedback they still miss branch interactions that the hidden tests exercise.

\subsection{Comparison with Refactoring Evaluation}

\textsc{MergeConflictBench} and \textsc{RefactorBench} share the hidden-test methodology but differ in what constitutes correctness:

\begin{table}[ht]
\centering
\caption{Comparison of behavioral evaluation in refactoring vs.\ merge conflict resolution.}
\label{tab:comparison}
\small
\begin{tabularx}{\textwidth}{lXX}
\toprule
\textbf{Dimension} & \textbf{RefactorBench} & \textbf{MergeConflictBench} \\
\midrule
Task & Decompose one file into modules & Integrate two divergent change sets \\
Correctness criterion & Preserve all existing behavior & Preserve base + ours + theirs behavior \\
Agent input & Single source file & Conflicted files with markers \\
Structural change & Expected (new files created) & Not expected (same files, markers removed) \\
Failure mode & Dropped behavior from refactoring & Dropped behavior from one branch \\
\bottomrule
\end{tabularx}
\end{table}

A key difference is that merge conflict resolution does not typically change the file structure---the output should be the same files as the input, but with markers replaced by resolved code. This makes the task more constrained than refactoring (where the agent has creative freedom over module boundaries) but also more demanding in terms of semantic understanding (where the agent must synthesize two competing implementations).

\subsection{Future Work}

\begin{itemize}[leftmargin=2em]
\item \textbf{Tool-using agent evaluation.} The reported model matrix uses a deliberately small one-shot block-replacement harness. Future runs should evaluate full coding agents with file editing, test execution on public smoke tests, iterative repair, and provider-specific APIs such as OpenAI's Responses API for models that are not exposed through chat completions.

\item \textbf{Language and framework diversity.} The current corpus is entirely JavaScript/JSX from React-family applications. Extending to Python, Go, Rust, and non-React frontend/backend frameworks would test generalization, as merge conflict patterns differ across languages and ecosystems (e.g., import conflicts in Python vs.\ TypeScript, ownership conflicts in Rust, component-state conflicts in React vs.\ template-driven frameworks).

\item \textbf{Automated test generation.} Manually authoring hidden test suites is the primary bottleneck. Exploring automated or semi-automated test generation from the three-way diff---potentially using an LLM to draft tests from the \texttt{preservedBehaviors} contract, with human review---could scale the benchmark significantly.

\item \textbf{Larger corpus via open-source repositories.} Extracting merge conflicts from popular open-source repositories with existing test suites could dramatically expand the corpus. If a repository's CI tests already cover the behavior introduced by each branch, those tests could serve as the hidden holdout suite with minimal additional authoring.

\item \textbf{Resolution strategy analysis.} Classifying agent resolutions by strategy (accept-ours, accept-theirs, accept-both, custom synthesis) and correlating with correctness would reveal which conflict types agents handle well and which they struggle with.
\end{itemize}

% ============================================================
\section{Conclusion}
\label{sec:conclusion}
% ============================================================

We introduced \textsc{MergeConflictBench}, a benchmark for evaluating LLM agents on merge conflict resolution using hidden behavioral tests. The benchmark extends the principle established by \textsc{RefactorBench}---behavioral preservation is a functional property best verified by functional tests---from single-branch refactoring to multi-branch integration.

The corpus consists of 86 real merge conflict cases from production React-family JavaScript/JSX code, spanning 288 conflicted files with 1{,}063 conflict blocks and 2{,}540 executable hidden tests across complexity tiers from 2-block import conflicts to 46-block fleet management workflows. All 86 cases have complete hidden test suites covering behaviors from all three code lineages (base, ours, theirs), ensuring that a correct resolution must genuinely integrate both branches' contributions rather than favoring one side.

Release-time validation confirms that the checked-in reference resolutions pass all hidden tests. Deterministic merge-editor baselines are substantially weaker: accept-ours, accept-theirs, and accept-both pass 17, 10, and 48 of 86 cases respectively. A one-shot LLM model matrix passes between 25 and 36 cases across five models, with Claude Opus 4.6 strongest in this harness at 36 of 86. These results show that the corpus requires more than conflict-marker removal or one-sided selection; a resolver must make semantic integration decisions, and current one-shot model calls still miss many of them.

Merge conflict resolution is a ubiquitous task in software development and a natural target for LLM automation. \textsc{MergeConflictBench} provides the evaluation infrastructure needed to develop and compare merge resolution agents with confidence that measured improvements correspond to genuine gains in semantic correctness.

% ============================================================
% References
% ============================================================

\bibliographystyle{plainnat}

\begin{thebibliography}{14}

\bibitem[Apel et~al.(2011)]{apel2011semistructured}
S.~Apel, O.~Lessening, and C.~Lengauer.
\newblock Structured merge revisited.
\newblock In \emph{Proceedings of the ACM SIGSOFT Symposium on the Foundations of Software Engineering (FSE)}, pages 190--201, 2011.

\bibitem[Chen(2025)]{chen2025refactorbench}
D.~Chen.
\newblock {RefactorBench}: Evaluating {LLM} agents on behavior-preserving code decomposition.
\newblock Technical report, Anything, Inc., 2025.

\bibitem[Dinella et~al.(2022)]{dinella2022deepmerge}
E.~Dinella, T.~Mytkowicz, A.~Svyatkovskiy, C.~Bird, and S.~Lahiri.
\newblock Deepmerge: Learning to merge programs.
\newblock In \emph{Proceedings of the ACM SIGSOFT Symposium on the Foundations of Software Engineering (FSE)}, 2022.

\bibitem[Gauthier(2024a)]{gauthier2024aider}
P.~Gauthier.
\newblock Aider refactoring benchmark.
\newblock \url{https://aider.chat/docs/leaderboards/refactor.html}, 2024.

\bibitem[Gauthier(2024b)]{gauthier2024polyglot}
P.~Gauthier.
\newblock Aider polyglot benchmark.
\newblock \url{https://aider.chat/2024/12/21/polyglot.html}, 2024.

\bibitem[Jimenez et~al.(2024)]{jimenez2024swebench}
C.~E. Jimenez, J.~Yang, A.~Wettig, S.~Yao, K.~Pei, O.~Press, and K.~Narasimhan.
\newblock {SWE-bench}: Can language models resolve real-world {GitHub} issues?
\newblock In \emph{Proceedings of the International Conference on Learning Representations (ICLR)}, 2024.

\bibitem[Sousa et~al.(2018)]{sousa2018verified}
M.~Sousa, I.~Dillig, and S.~Lahiri.
\newblock Verified three-way program merge.
\newblock \emph{Proceedings of the ACM on Programming Languages}, 2(OOPSLA):1--29, 2018.

\bibitem[Svyatkovskiy et~al.(2022)]{svyatkovskiy2022mergebert}
A.~Svyatkovskiy, T.~Mytkowicz, E.~Dinella, C.~Bird, and S.~Lahiri.
\newblock {MergeBERT}: Program merge conflict resolution via neural multi-task learning.
\newblock In \emph{Proceedings of the ACM Joint European Software Engineering Conference and Symposium on the Foundations of Software Engineering (ESEC/FSE)}, 2022.

\bibitem[Wang et~al.(2024)]{wang2024openhands}
X.~Wang, B.~Li, Y.~Song, F.~F. Xu, X.~Tang, M.~Zhuge, J.~Pan, Y.~Song, B.~Li, J.~Singh, H.~H. Tran, F.~Li, R.~Ma, M.~Zhang, B.~Yu, J.~Bisk, C.~Xiong, T.~Yu, G.~Neubig, and Y.~Li.
\newblock Open{H}ands: An open platform for {AI} software developers as generalist agents.
\newblock \emph{arXiv preprint arXiv:2407.16741}, 2024.

\bibitem[Yang et~al.(2024)]{yang2024sweagent}
J.~Yang, C.~E. Jimenez, A.~Wettig, K.~Lieret, S.~Yao, K.~Narasimhan, and O.~Press.
\newblock {SWE}-agent: Agent-computer interfaces enable automated software engineering.
\newblock \emph{arXiv preprint arXiv:2405.15793}, 2024.

\bibitem[Zhang et~al.(2024)]{zhang2024mergegen}
Y.~Zhang, H.~Yu, Z.~Sun, and J.~Dong.
\newblock {MergeGen}: A retrieval-augmented approach for automatic merge conflict resolution.
\newblock \emph{arXiv preprint arXiv:2403.07766}, 2024.

\end{thebibliography}

\end{document}
